\documentclass[11pt,a4paper]{article}
\usepackage[T1]{fontenc}
\usepackage{lmodern}
\usepackage[margin=25mm]{geometry}
\usepackage{amsmath,amssymb,amsthm}
\usepackage[hidelinks]{hyperref}
\usepackage{microtype}
\newcommand{\Q}{\mathbb Q}
\newcommand{\R}{\mathbb R}
\newcommand{\N}{\mathbb N}
\title{Divergent mean heights on a nonpreperiodic quadratic orbit}
\author{Submission for TLMC problem 00000002689}
\date{24 October 2026}
\begin{document}
\maketitle

\paragraph{Scope.}
The conjecture asserts that the mean naive height of an orbit outside the
preperiodic set converges, with a proposed canonical-height expression and
rate. We use the ordinary arithmetic mean over the first $n$ iterates,
without degree normalization. For the degree-two polynomial $f(z)=z^2$
over $\Q$, the orbit of $2$ disproves this necessary convergence assertion.
Both standard conventions for naive height, logarithmic and
multiplicative, give means tending to $+\infty$.
We do not specify the source's undefined Fourier coefficients or
orbit-frequency measure, and make no claim against a different,
degree-normalized averaging statement.

\paragraph{Actual heights and iterates.}
For a rational number $a/b$ in lowest terms with $b>0$, the standard
multiplicative and logarithmic Weil heights are
\[
H(a/b)=\max\{|a|,b\},\qquad h(a/b)=\log H(a/b).
\]
The formal proof uses Mathlib's native heights defined from the
archimedean and nonarchimedean absolute values. The displayed rational
formulas are proved consequences of those definitions \cite{height}.
In particular, no orbit height is assigned by definition.

Take $f(z)=z^2$ and $x_0=2$. The actual iterates satisfy
\[
x_k=f^k(2)=2^{2^k},\qquad
H(x_k)=2^{2^k},\qquad h(x_k)=2^k\log2.
\]
The sequence $x_k$ is strictly increasing. Its forward orbit is therefore
infinite, so $2$ is not preperiodic. The formalization constructs the
actual polynomial $X^2$, proves its degree is two and its evaluation is
$f$, and defines preperiodicity by finiteness of the actual forward orbit.

\paragraph{The canonical height is finite.}
The canonical height of a degree-$d>1$ map is obtained from the normalized
height limit \cite{demarco}. Here, for every $x\in\Q$ and every $k\in\N$,
\[
\frac{h(f^k(x))}{2^k}=h(x).
\]
Consequently the genuine normalized limit exists and
\[
\widehat h_f(x)=\lim_{k\to\infty}\frac{h(f^k(x))}{2^k}=h(x),
\qquad \widehat h_f(2)=\log2>0.
\]
In Lean the limit is defined by \texttt{Filter.limUnder}; convergence is
proved before its value is identified. Thus the finite canonical height
comes from its defining limit, not from a numerical placeholder.

\paragraph{The unnormalized mean diverges.}
For $n\ge1$, let
\[
A_n=\frac1n\sum_{k=0}^{n-1}h(x_k).
\]
The actual finite geometric sum gives
\[
A_n=\frac{2^n-1}{n}\log2.
\]
For completeness, an elementary lower bound suffices to establish
divergence without an asymptotic approximation. The inequality
$2^k\ge k+1$ implies
\[
A_n\ge\frac{\log2}{n}\sum_{k=0}^{n-1}(k+1)
=\frac{n+1}{2}\log2\longrightarrow+\infty.
\]
Therefore $A_n$ cannot converge to any finite real number. In particular,
it cannot have the finite limit asserted in the conjecture, and hence
cannot converge to such a limit at rate $1/\log n$.

If ``naive height'' instead means $H$, the conclusion is unchanged.
Because $h(x)=\log H(x)\le H(x)$ and $H(x)\ge1$, the arithmetic means
\[
B_n=\frac1n\sum_{k=0}^{n-1}H(x_k)
\]
satisfy $B_n\ge A_n$ and also tend to $+\infty$. The formal conclusion
excludes convergence to every real number for both sequences.

\paragraph{Formal scope and verification.}
The complete proof is in \texttt{lean/Main.lean}, culminating in
\texttt{TLMC2689.counterexample}. The statement connects the actual
degree-two polynomial, all its iterates, the nonpreperiodic orbit, the
native heights, the canonical-height limit and both divergent means.
All convergence assertions concern $n\to\infty$; the lower-bound argument
uses $n\ge1$ throughout. No assertion depends on a zero-denominator value.

The original bilingual conjecture is preserved in \texttt{SOURCE.md}.
The project pins Lean 4.33.0 and Mathlib commit
\begin{center}\small
\texttt{db584cd6d46c92f209a44c0f1c829460d327499d}.
\end{center}
The accompanying evidence records the exact build and exhaustive
namespace axiom audit, a fresh kernel replay, rendered-PDF review and
file hashes. The disproof concerns the ordinary, unnormalized mean only;
it does not dispute canonical-height theory or arithmetic
equidistribution under their hypotheses.

\begin{thebibliography}{2}
\bibitem{height}
Mathlib, \emph{Heights on number fields},
\href{https://github.com/leanprover-community/mathlib4/blob/db584cd6d46c92f209a44c0f1c829460d327499d/Mathlib/NumberTheory/Height/NumberField.lean}{pinned source}, rational-height formulas.
\bibitem{demarco}
L.~DeMarco, \emph{Arithmetic Dynamics and Intersection Problems},
Arizona Winter School 2023, \S1.4, p.~6,
\href{https://people.math.harvard.edu/~demarco/Math274/AWS_LectureNotes.pdf}{author's lecture notes}.
\end{thebibliography}
\end{document}
